% docs/slides/preamble.tex -- 五份 deck 共享
\documentclass[aspectratio=169,11pt]{beamer}

% ==== 中文与字体 ====
\usepackage[UTF8,fontset=none]{ctex}
\usepackage{fontspec}

\IfFontExistsTF{PingFang SC}{%
  \setCJKmainfont{PingFang SC}[BoldFont=PingFang SC,ItalicFont=PingFang SC,AutoFakeSlant=0.15]
  \setCJKsansfont{PingFang SC}[ItalicFont=PingFang SC,AutoFakeSlant=0.15]
  \setCJKmonofont{PingFang SC}
}{%
  \IfFontExistsTF{Source Han Sans SC}{%
    \setCJKmainfont{Source Han Sans SC}
    \setCJKsansfont{Source Han Sans SC}
  }{%
    \IfFontExistsTF{Noto Sans CJK SC}{%
      \setCJKmainfont{Noto Sans CJK SC}
      \setCJKsansfont{Noto Sans CJK SC}
    }{%
      \setCJKmainfont{Heiti SC}
      \setCJKsansfont{Heiti SC}
    }
  }
}

\IfFontExistsTF{Helvetica Neue}{\setmainfont{Helvetica Neue}}{%
  \IfFontExistsTF{Inter}{\setmainfont{Inter}}{}}
\IfFontExistsTF{JetBrains Mono}{\setmonofont{JetBrains Mono}[Scale=0.85]}%
  {\setmonofont{Menlo}[Scale=0.85]}

% ==== 颜色 ====
\usepackage{xcolor}
\definecolor{cMain}{HTML}{1F4E79}
\definecolor{cAccent}{HTML}{E07A1F}
\definecolor{cGray}{HTML}{595959}
\definecolor{cMute}{HTML}{F4F4F4}
\definecolor{cOK}{HTML}{2E7D32}
\definecolor{cBad}{HTML}{B23A3A}
\definecolor{cWarn}{HTML}{C8A300}

% ==== 主题 ====
\usepackage{tcolorbox}
\usetheme{metropolis}
\metroset{progressbar=frametitle,numbering=fraction,block=fill}
\setbeamercolor{frametitle}{bg=cMain,fg=white}
\setbeamercolor{progress bar}{fg=cAccent,bg=cMute}
\setbeamercolor{alerted text}{fg=cAccent}
\setbeamercolor{title}{fg=cMain}

% metropolis 默认尝试 Fira Sans，未安装时 fallback 到 Latin Modern Sans。
% 这里在主题加载后显式覆盖，强制使用 macOS 上一定存在的 Helvetica Neue，
% 跟中文 PingFang 风格一致；找不到则按顺序尝试 Inter / Arial。
\IfFontExistsTF{Helvetica Neue}{%
  \setsansfont{Helvetica Neue}%
}{%
  \IfFontExistsTF{Inter}{\setsansfont{Inter}}{%
    \IfFontExistsTF{Arial}{\setsansfont{Arial}}{}}%
}

% ==== TikZ ====
\usepackage{tikz}
\usetikzlibrary{positioning,arrows.meta,fit,shapes.geometric,calc,%
  decorations.pathmorphing,backgrounds}

\tikzset{
  node-box/.style={draw=cMain,thick,rounded corners=2pt,
    minimum height=8mm,minimum width=18mm,align=center,font=\small},
  node-state/.style={draw=cMain,thick,circle,minimum size=14mm,
    align=center,font=\small},
  node-ok/.style={draw=cOK,thick,fill=cOK!10},
  node-warn/.style={draw=cWarn,thick,fill=cWarn!15},
  node-bad/.style={draw=cBad,thick,fill=cBad!10},
  % 色盲友好的状态机三态：蓝 (正常/Closed) / 橙 (告警/Open) / 灰 (中间/HalfOpen)
  % 避免 red-green 二元区分，改用 hue + lightness 双通道。
  node-state-normal/.style={draw=cMain,thick,fill=cMain!12},
  node-state-alert/.style={draw=cAccent,thick,fill=cAccent!18},
  node-state-transition/.style={draw=cGray,thick,fill=cGray!18},
  arrow/.style={-{Stealth[length=2mm]},thick,draw=cMain},
  arrow-async/.style={-{Stealth[length=2mm]},thick,dashed,draw=cGray},
  arrow-bad/.style={-{Stealth[length=2mm]},thick,draw=cBad},
  arrow-alert/.style={-{Stealth[length=2mm]},thick,draw=cAccent},
  lifeline/.style={dashed,thick,draw=cGray},
}

% ==== 代码 ====
\usepackage{listings}
\lstdefinelanguage{Go}{
  morekeywords={break,case,chan,const,continue,default,defer,else,
    fallthrough,for,func,go,goto,if,import,interface,map,package,range,
    return,select,struct,switch,type,var,nil,true,false,iota},
  morekeywords=[2]{string,int,int64,uint,uint64,bool,byte,error,context},
  sensitive=true,
  morecomment=[l]//,morecomment=[s]{/*}{*/},
  morestring=[b]",morestring=[b]`,
}
\lstdefinelanguage{LuaRedis}{
  morekeywords={and,break,do,else,elseif,end,false,for,function,if,in,
    local,nil,not,or,repeat,return,then,true,until,while},
  morekeywords=[2]{redis,call,KEYS,ARGV,tonumber,HGET,HSET,GET,SET,
    DECR,DECRBY,INCR,INCRBY,EXPIRE,PEXPIRE,ZADD,ZCARD,ZREMRANGEBYSCORE,
    EVAL,EVALSHA},
  sensitive=true,
  morecomment=[l]{--},
  morestring=[b]",morestring=[b]',
}

\lstdefinestyle{base}{
  basicstyle=\ttfamily\scriptsize,
  numbers=left,numberstyle=\tiny\color{cGray},
  keywordstyle=\color{cMain}\bfseries,
  keywordstyle=[2]\color{cAccent},
  stringstyle=\color{cOK},
  commentstyle=\color{cGray}\itshape,
  backgroundcolor=\color{cMute},
  frame=single,framerule=0pt,
  showstringspaces=false,
  breaklines=true,
  tabsize=2,
  columns=fullflexible,
}
\lstset{style=base}

% ==== 三线表与附录页码 ====
\usepackage{booktabs}
\usepackage{appendixnumberbeamer}

% ==== 页脚来源标注 ====
\newcommand{\srcnote}[1]{{\color{cGray}\scriptsize\itshape #1}}

% 关键论点框
\newcommand{\keypoint}[1]{%
  \begin{tcolorbox}[colback=cMute,colframe=cMain,boxrule=0.5pt,arc=2pt]%
  #1\end{tcolorbox}}


\title{库存模型的业务承诺}
\subtitle{两桶库存、Saga 取消与 30 分钟那条线}
\author{gomall · 业务侧 deck}
\date{}

\begin{document}

\maketitle

\begin{frame}{今日大纲}
\tableofcontents
\end{frame}

% ============================================================
\section{一笔订单背后，库存做了三件事}
% ============================================================

\begin{frame}{从用户视角看一次下单}
\begin{itemize}
  \item 用户点"立即购买"：期待\textbf{马上锁住}，不能被别人抢走。
  \item 用户没付钱跑了：期待"商品很快回到货架"，否则商家亏。
  \item 用户付完钱：期待"这件东西就是我的"，不能再次卖给别人。
  \item 这三件事，对系统都是\textbf{同一个库存数字}，但要给出三种不同承诺。
\end{itemize}
\srcnote{service/order.go:40--102 OrderCreate 主路径}
\end{frame}

\begin{frame}{两桶库存：可售与已锁定}
\begin{tikzpicture}[node distance=8mm and 14mm]
  \node[node-box,node-ok] (av) {available\\\small 可售};
  \node[node-box,node-warn,right=of av] (rv) {reserved\\\small 已锁定};
  \node[node-box,below=of av] (db) {DB product.num\\\small 真实水位};
  \draw[arrow] (av.north east) to[bend left=20] node[above,font=\tiny]{下单 reserve} (rv.north west);
  \draw[arrow] (rv.south west) to[bend left=20] node[below,font=\tiny]{取消 release} (av.south east);
  \draw[arrow-async] (db) -- node[right,font=\tiny]{启动 seed} (av);
  \draw[arrow-bad] (rv.east) -- ++(1.2,0) node[right,font=\tiny]{支付 commit\\\tiny 真扣减};
\end{tikzpicture}
\vspace{2mm}
\begin{itemize}
  \item \textbf{available} = 给前端看的"还能下几单"
  \item \textbf{reserved} = 已经被人占着、但还没付钱的份额
  \item 二者之和始终 $\le$ 商品总水位
\end{itemize}
\srcnote{repository/cache/inventory.go:11--22 两桶 Key}
\end{frame}

\begin{frame}{业务为什么要分两桶}
\begin{itemize}
  \item 只有一桶时只有两个选择：
  \begin{itemize}
    \item 下单立刻扣 DB → 用户跑了得回写，回写期间别人看到的库存不准。
    \item 下单不扣 DB → 100 个人抢 5 件，付款时 95 人失败，比"卖光"更糟。
  \end{itemize}
  \item 分两桶：available 给"用户视图"，reserved 给"运营对账"。
  \item 30 分钟内付钱：reserved $\to$ 真扣；不付钱：reserved $\to$ available 回滚。
  \item 业务上能说清楚：\textbf{已售 = 已支付，已锁 = 占座未结账}。
\end{itemize}
\srcnote{repository/cache/inventory.go:29--66 三段 Lua 脚本对应三种业务动作}
\end{frame}

% ============================================================
\section{30 分钟自动关单的来历}
% ============================================================

\begin{frame}{30 分钟是怎么定出来的}
\begin{tabular}{@{}llp{4.5cm}@{}}
\toprule
方案 & 用户体验 & 业务后果 \\
\midrule
5 分钟   & 来不及切支付方式      & 商家投诉"刚下单就被取消" \\
15 分钟  & 多数人够用           & 高峰期仍有 5\%-8\% 误关 \\
\textbf{30 分钟}  & 含切卡 / 收验证码冗余 & 大促锁单可控 \\
60 分钟  & 用户基本不抱怨        & 双 11 锁 1 小时谁也抢不到 \\
24 小时  & 客服零投诉           & 库存周转崩盘 \\
\bottomrule
\end{tabular}
\vspace{2mm}
\small 30 分钟是\textbf{商家投诉}与\textbf{库存周转}两条曲线的交点，行业普遍落在 15--45min。
\srcnote{repository/rabbitmq/order\_delay.go:22 OrderCancelDelay = 30 * time.Minute}
\end{frame}

\begin{frame}{30min 这条线触发了哪些动作}
\begin{tikzpicture}[node distance=10mm and 14mm]
  \node[node-box] (u) {用户下单};
  \node[node-box,right=of u,node-ok] (r) {Redis reserve};
  \node[node-box,right=of r] (db) {DB 写 order\\(UnPaid)};
  \node[node-box,below=of db,node-warn] (mq) {RMQ delay\\30min TTL};
  \node[node-box,below=of r,node-warn] (cron) {Cron\\每 5min 兜底};
  \node[node-box,left=of cron,node-state-alert] (cn) {Cancel\\release reserve};
  \draw[arrow] (u) -- (r); \draw[arrow] (r) -- (db); \draw[arrow] (db) -- (mq);
  \draw[arrow-alert] (mq) -- node[above,font=\tiny]{TTL 到} (cron);
  \draw[arrow] (cron) -- (cn);
  \draw[arrow-async] (db.south west) to[bend right=15] node[left,font=\tiny]{MQ 挂} (cn.north east);
\end{tikzpicture}
\vspace{2mm}
\small 30min 不是一个定时器，是\textbf{两条并行的兜底链}：MQ 延迟队列 + Cron 扫表。
\srcnote{repository/rabbitmq/order\_delay.go:59--72 发布；initialize/cron.go:13--28 Cron}
\end{frame}

\begin{frame}{兜底链为什么不能只留一条}
\begin{itemize}
  \item \textbf{只 MQ}：RabbitMQ 进程挂 / 队列丢消息 $\to$ 订单永远停在 UnPaid。
  \item \textbf{只 Cron}：5min 间隔 + 单批 100 单上限，库存释放 lag 大。
  \item gomall 选择两条都留：
  \begin{itemize}
    \item MQ 是\textbf{主路径}，正常情况 30min 准点触发。
    \item Cron 每 5 分钟扫一次 \texttt{type=UnPaid AND created\_at \le now-15min} 的订单兜底。
    \item \textbf{业务保证}：哪怕 MQ 彻底挂掉，库存最多多锁 35 分钟，不会泄漏到第二天。
  \end{itemize}
\end{itemize}
\srcnote{service/order\_task.go:14--48 Cron；repository/db/dao/order.go:162--170 扫表}
\end{frame}

% ============================================================
\section{超卖 vs 误锁：哪个客诉更难处理}
% ============================================================

\begin{frame}{这道选择题没有"两不"的选项}
\begin{tikzpicture}[node distance=10mm and 16mm]
  \node[node-state,node-state-alert] (over)  {超卖};
  \node[node-state,node-state-transition,right=of over]   (mid)   {临界点};
  \node[node-state,node-state-normal,right=of mid]   (miss)  {误锁};
  \draw[arrow-alert] (over) -- node[above,font=\tiny]{补差价 / 道歉} (mid);
  \draw[arrow]       (mid)  -- node[above,font=\tiny]{用户走开} (miss);
\end{tikzpicture}
\vspace{2mm}
\begin{tabular}{@{}lll@{}}
\toprule
事故 & 用户感知 & 处置成本 \\
\midrule
\textbf{超卖} & "已下单 + 已付款"，最后没货 & 真金白银赔 + 商誉 \\
\textbf{误锁} & 货架上写"售罄"，没下单           & 用户走开，无合同纠纷 \\
\bottomrule
\end{tabular}
\vspace{2mm}
\small 已下单是合同。系统的默认偏好必须是\textbf{宁可误锁，不可超卖}。
\srcnote{repository/cache/inventory.go:32--44 reserveScript 先判 avail 再 DECRBY}
\end{frame}

\begin{frame}{怎么把"宁可误锁"翻译成代码}
\begin{itemize}
  \item Lua 脚本第一行就是 \texttt{if avail $<$ need then return -1}：
  库存够才扣，不够直接拒绝，没有"先扣了再说"的中间态。
  \item Redis 单线程 + Lua 原子执行：500 并发抢 100 件券，单元测试\textbf{零超发}。
  \item DB 那一桶（\texttt{product.num}）只在\textbf{支付落地}和\textbf{Cron 回滚}两条路径上动。
  \item 业务承诺：\textbf{任何瞬间，已售 + 已锁 $\le$ 总水位}，这条不变量永不破。
\end{itemize}
\srcnote{stressTest/REPORT.md \S 单元测试 TestCoupon\_AtomicClaimNoOversell 500 并发 / 100 张 / 零超发}
\end{frame}

\begin{frame}{两个错误码，两套客服话术}
\begin{tabular}{@{}lll@{}}
\toprule
触发 & 业务现象 & 客服话术 \\
\midrule
ErrStockInsufficient & available $<$ need     & "已售罄，下次抢" \\
ErrStockNotInit       & available 桶不存在     & "商品上下架中" \\
\bottomrule
\end{tabular}
\vspace{2mm}
\begin{itemize}
  \item 第二个不能说"库存系统异常" —— 用户听不懂、客服没法答。
  \item 内部按 \texttt{ErrStockNotInit} 报 SRE，外部统一是"商品上下架中"。
  \item 不向用户暴露 Redis / Lua / 桶这些词。
\end{itemize}
\srcnote{repository/cache/inventory.go:25--27 两个 sentinel error}
\end{frame}

% ============================================================
\section{Redis 扣成功、DB 落库失败的那一秒}
% ============================================================

\begin{frame}{下单的故障窗口}
\begin{tikzpicture}[node distance=8mm and 14mm]
  \node[node-box] (s0) {收到请求};
  \node[node-box,right=of s0,node-ok] (s1) {Redis reserve};
  \node[node-box,right=of s1,node-warn] (s2) {DB Tx};
  \node[node-box,right=of s2,node-ok] (s3) {RMQ publish};
  \node[node-box,below=of s2,node-state-alert] (x) {进程 crash\\Redis 已扣 / DB 未落};
  \draw[arrow] (s0) -- (s1); \draw[arrow] (s1) -- (s2); \draw[arrow] (s2) -- (s3);
  \draw[arrow-bad] (s1.south) to[bend right=10] (x.north west);
  \draw[arrow-bad] (s2.south) -- (x.north);
\end{tikzpicture}
\vspace{2mm}
\small 风险窗口很小，但只要并发 $\times$ 时长 $>$ 0，就会发生。\\
gomall 的处理：\textbf{Tx 失败立刻 release}，不留给 Cron。
\srcnote{service/order.go:79--85 事务失败显式 ReleaseReservation}
\end{frame}

\begin{frame}{Saga 取消：把"已扣"还回去}
\begin{tikzpicture}[node distance=8mm and 12mm]
  \node[node-state,node-state-normal] (un) {UnPaid};
  \node[node-state,node-state-transition,right=of un] (cl) {Closing};
  \node[node-state,node-state-alert,right=of cl] (cn) {Cancelled};
  \draw[arrow] (un) -- node[above,font=\tiny]{30min TTL 到} (cl);
  \draw[arrow] (cl) -- node[above,font=\tiny]{release + outbox} (cn);
  \draw[arrow-async] (cn.south) to[bend left=30] node[below,font=\tiny]{二次调用 no-op} (cn.north);
\end{tikzpicture}
\vspace{2mm}
\begin{itemize}
  \item \textbf{CloseOrderWithCheck} 用 WHERE type=UnPaid 自带幂等：状态机已经流过的不再变。
  \item release reserve 也是幂等的：\texttt{reserved $<$ n} 直接报错并跳过，不会出现负数。
  \item Saga 的本质：\textbf{业务两段，回滚两段，每段都能重放}。
\end{itemize}
\srcnote{service/order\_cancel.go:19--63；repository/db/dao/order.go:185--195}
\end{frame}

\begin{frame}{为什么不直接 DB 事务包住 Redis}
\begin{itemize}
  \item Redis 没有跨进程事务，参与不了 DB 的 2PC。
  \item 强行加分布式事务（XA / TCC）：性能拉胯、运维复杂、得不偿失。
  \item gomall 选择 Saga + 兜底：
  \begin{itemize}
    \item 正常路径\textbf{先 Redis 后 DB}，因为 Redis 失败更廉价。
    \item 异常路径\textbf{显式回滚 Redis} + Cron 二次兜底。
    \item \textbf{业务承诺}：\textbf{最终一致}，不追求"每秒一致"。
  \end{itemize}
  \item 这一选择对应技术 deck 04：Outbox + Saga 的最终一致性。
\end{itemize}
\srcnote{service/order.go:57--98 reserve $\to$ tx $\to$ publish 三段}
\end{frame}

% ============================================================
\section{库存 leak 的几种典型场景}
% ============================================================

\begin{frame}{leak 是什么、为什么不能忍}
\begin{itemize}
  \item \textbf{库存 leak}：reserved 桶里的数字回不到 available。
  \item 直接后果：商品看起来"无货"，但商家仓库里还有。
  \item 商家视角：\textbf{每分钟都在丢单}，看不到也修不了。
  \item 工程师视角：\textbf{一个隐形的钱漏}，问题积累一周才会被发现。
  \item 所以 leak 不能靠"发现了再修"，必须主动巡检。
\end{itemize}
\srcnote{repository/cache/inventory.go:128--138 GetStockSnapshot 对账接口}
\end{frame}

\begin{frame}{四种典型 leak 场景}
\begin{tabular}{@{}lll@{}}
\toprule
场景 & 谁的锅 & 兜底机制 \\
\midrule
进程 crash 在 reserve 后 / tx 前 & 偶发 & Cron 扫 UnPaid \\
DB tx 失败但 release 也失败       & Redis 抖动 & Cron 兜底 \\
RMQ 消费失败 (nack 进 DLX)        & 消费者 bug & Cron 兜底 \\
Cron 自身没跑到 (调度挂)          & 运维 & 监控告警 \\
\bottomrule
\end{tabular}
\vspace{2mm}
\small Cron 是\textbf{最后一道线}，所以扫表条件是 \texttt{created\_at \le now-15min} 而不是 30min ——\\
\textbf{15min 早于 MQ 的 30min}，保证 MQ 正常时 Cron 不会抢先误关。
\srcnote{service/order\_task.go:16 GetTimeoutOrders(15, 100)}
\end{frame}

\begin{frame}{对账：每天都要回答的那一句话}
\begin{itemize}
  \item 凌晨低峰跑对账脚本（未实现，运维 backlog 第 1 项）：
  \begin{itemize}
    \item 对每个 productID 读 \texttt{GetStockSnapshot} $\to$ (available, reserved)
    \item 再统计 \texttt{COUNT(order WHERE type=UnPaid AND product\_id=?)}
    \item \textbf{期望}：\texttt{reserved == SUM(unpaid.num)}
    \item 偏差 $>$ 1 件 / $>$ 10min：钉钉告警 + 自动修复或人工介入
  \end{itemize}
  \item 这一步\textbf{业务必须能跑}，否则 30min 这条线名存实亡。
  \item 类似支付系统的"日切对账"，跑不出 $\to$ 不许结账。
\end{itemize}
\srcnote{repository/cache/inventory.go:127--138 给对账脚本的入口}
\end{frame}

% ============================================================
\section{四个角色看同一件事}
% ============================================================

\begin{frame}{用户 / 运营 / 商家 / SRE 各看到什么}
\begin{tabular}{@{}llp{5.5cm}@{}}
\toprule
角色 & 关心的事 & 看哪个指标 \\
\midrule
用户   & 我抢到了吗               & 下单成功 + 30min 倒计时 \\
运营   & 库存有没有泄漏           & reserved / available 比值 \\
商家   & 会不会被恶意锁单         & 同一用户的待支付订单数 \\
SRE   & Redis 和 DB 一致吗      & 对账脚本偏差秒数 \\
\bottomrule
\end{tabular}
\vspace{3mm}
\small 同一份代码，四个口径。每个口径背后都有\textbf{独立 SLO}：\\
用户侧 99.9\% 一次下单成功；SRE 侧 99.95\% 一致性。
\srcnote{routes/router.go:77 orders/create 挂幂等中间件}
\end{frame}

\begin{frame}{恶意锁单：商家最怕的事}
\begin{itemize}
  \item 攻击者用脚本反复下单、不付款，把热门商品的 reserved 桶刷满。
  \item 真用户看不到货，攻击者拖到下一批补货前继续刷。
  \item 防线：
  \begin{itemize}
    \item 路由层：\texttt{/orders/create} 挂 Idempotency 中间件（同 Key 只生效一次）。
    \item 业务层（建议补）：单用户待支付订单数 $\le$ 5（参考 deck 10 限券方案）。
    \item 风控层：脚本特征识别（不在 gomall 范围内）。
  \end{itemize}
  \item 业务承诺：\textbf{30min 内能锁的总量}，不是无界的。
\end{itemize}
\srcnote{routes/router.go:77 middleware.Idempotency 接入}
\end{frame}

% ============================================================
\section{业务承诺一张总表}
% ============================================================

\begin{frame}{压测数据替业务说话}
\begin{tabular}{@{}lrl@{}}
\toprule
场景 & 实测 & 业务含义 \\
\midrule
500 并发抢 100 张      & 零超发     & 不变量守住 \\
50 VU 同 Key 重放 15s & 755{,}033 次 / 1 单 & 幂等中间件接住恶意锁单 \\
/orders/create RPS  & 50{,}319   & 单机够撑双 11 中等品类 \\
p95 / max           & 2.33ms / 73ms & 远低于 P0 SLO 500ms \\
\bottomrule
\end{tabular}
\vspace{2mm}
\small 这些不是"性能数字"，是\textbf{业务承诺的实测证据}。\\
压测脚本和数据都在 \texttt{stressTest/} 下，PR 评审可复核。
\srcnote{stressTest/REPORT.md 链路压测表 第 5 行 + \S 单元测试}
\end{frame}

\begin{frame}{业务承诺总表}
\begin{tabular}{@{}lll@{}}
\toprule
承诺 & 实现 & 触发条件 \\
\midrule
\textbf{不超卖}      & Redis Lua 原子扣 & 100\% 时刻 \\
\textbf{30min 释放} & MQ TTL 30min     & 用户未付款 \\
\textbf{兜底释放}    & Cron 5min / 15min 窗口 & MQ 挂或消费失败 \\
\textbf{支付即扣}    & CommitReservation & 支付回调成功 \\
\textbf{幂等关单}    & WHERE type=UnPaid & 多次取消同一单 \\
\textbf{显式回滚}    & service/order.go release & DB tx 失败 \\
\bottomrule
\end{tabular}
\vspace{2mm}
\small 六条承诺缺一不可。任意一条破，业务侧就要回退到"人工对账"。
\end{frame}

\begin{frame}{回到那条 30 分钟的线}
\begin{itemize}
  \item 30 分钟是\textbf{用户友好}与\textbf{库存周转}的平衡。
  \item 它不是 SLA、不是延迟、不是性能指标 —— 它是\textbf{业务的合同条款}。
  \item 改 30 分钟，要走的不是一个 PR，是一次跨部门的产品决策：
  \begin{itemize}
    \item 商家：会不会涨锁单时长？
    \item 运营：节日要不要临时缩短？
    \item 风控：黑产能不能利用？
  \end{itemize}
\end{itemize}
\keypoint{库存治理的核心，不是 Redis Lua 怎么写，是\textbf{这六条承诺怎么落到代码里且能被压测复核}。}
\end{frame}

% ============================================================
\appendix
% ============================================================

\begin{frame}{面试 Q\&A}
\textbf{Q1}：为什么 30min，不是 1h？\\
\hspace{4mm}\small A：商家投诉曲线 + 双 11 锁单周转。15min 误关高、1h 大促锁不动，30min 是工业界经验值。

\vspace{2mm}
\textbf{Q2}：Redis 扣了但 DB 没落库怎么办？\\
\hspace{4mm}\small A：service/order.go 在 tx 失败分支显式 release；外加 Cron 兜底窗口 15min。两层。

\vspace{2mm}
\textbf{Q3}：库存 leak 怎么发现？\\
\hspace{4mm}\small A：GetStockSnapshot 出对账脚本，比对 reserved 与 SUM(UnPaid.num)，偏差超过阈值告警。

\vspace{2mm}
\textbf{Q4}：为什么不上分布式事务？\\
\hspace{4mm}\small A：Redis 不参与 2PC；Saga + 兜底成本更低，已经满足"最终一致"业务承诺。
\end{frame}

\begin{frame}{业务侧总结}
\begin{itemize}
  \item 两桶库存翻译成业务语言 = \textbf{可售} / \textbf{已锁定}。
  \item 30 分钟是合同条款，不是定时器参数。
  \item 超卖必赔，误锁仅流失 —— 默认偏好\textbf{宁可误锁}。
  \item 一致性靠 Saga 与对账，而不是分布式事务。
  \item 每条业务承诺，都\textbf{要有压测或单测能复核}，否则只是口号。
\end{itemize}
\keypoint{对应技术 deck：02 防超发、04 Outbox + Saga。这份业务侧 deck 解释的是"为什么这么做"，不替代它们。}
\end{frame}

\begin{frame}{代码位置一览}
\begin{description}
  \item[两桶 Key]\srcnote{repository/cache/inventory.go:16--22}
  \item[reserve Lua]\srcnote{repository/cache/inventory.go:32--44}
  \item[commit Lua]\srcnote{repository/cache/inventory.go:48--55}
  \item[release Lua]\srcnote{repository/cache/inventory.go:58--66}
  \item[下单主路径]\srcnote{service/order.go:40--102}
  \item[Saga 取消]\srcnote{service/order\_cancel.go:19--63}
  \item[Cron 兜底]\srcnote{service/order\_task.go:14--48}
  \item[MQ 30min 延迟]\srcnote{repository/rabbitmq/order\_delay.go:22--72}
  \item[启动 seed]\srcnote{service/inventory/syncer.go:14--48}
  \item[Cron 调度]\srcnote{initialize/cron.go:13--39}
  \item[对账入口]\srcnote{repository/cache/inventory.go:127--138}
\end{description}
\end{frame}

\end{document}
